\section{Évaluation}
\redacteur{\acompleter[nom]}

\consigne{Section la plus importante pour le critère P3. Mesurez chaque critère de réussite défini en semaine 1 sur le jeu de test, comparez à la cible, puis analysez les erreurs avec des exemples concrets. Un résultat inférieur à la cible, bien analysé, vaut mieux qu'un bon chiffre non expliqué.}

\subsection{Résultats}

\begin{table}[H]\centering
\caption{Critères de réussite : cibles et résultats}\label{tab:resultats}
\begin{tabularx}{\linewidth}{Xlll}
\toprule
\textbf{Indicateur} & \textbf{Cible} & \textbf{Résultat} & \textbf{Atteint} \\
\midrule
\acompleter & \acompleter & \acompleter & oui / non \\
\acompleter & \acompleter & \acompleter & oui / non \\
\acompleter & \acompleter & \acompleter & oui / non \\
\bottomrule
\end{tabularx}
\end{table}

\subsection{Protocole}
\acompleter[comment les mesures ont été faites : jeu de test, nombre d'exécutions, matériel, modèle et paramètres]

\subsection{Analyse des erreurs}
\consigne{Classez les échecs par type (mauvaise recherche, hallucination, mauvais outil appelé, délai\ldots), donnez leur fréquence et au moins deux exemples commentés.}
\acompleter
